Rice leaf photos are taken at high resolution, but most recognition systems shrink them to a few
hundred pixels, and many reports do not say how the data were split.
We study a resolution-preserving alternative that needs only photo-level labels. Each photo is
brought to a common leaf scale and cut into 224-pixel tiles. A frozen DINOv3 encoder describes every tile once,
and a small gated-attention head, trained from the photo label alone, combines the descriptions and
produces a heat map. We compare it with four fine-tuned whole-image models, probes, simpler pooling rules and
ablations. The protocol groups near-duplicates and capture sessions, reports macro-F1 and balanced accuracy with
spread and bootstrap ranges, and adds a zero-shot test on field photos (Paddy Doctor).
On the session-blocked split the head reaches 77.9\,\% test macro-F1. It is clearly above two of the four
whole-image models at 224 px and above both whole-photo controls, and not clearly different from ConvNeXt-T at
448 px, which has the higher mean. A committee of the two is clearly better than either.
Zero-shot transfer to field photos is not shown: all methods score low.
Removing the most-attended tiles hurts far more than removing random tiles.
Augmenting before splitting raises the score by 11.5 to 17.2 points, which shows how much the split order matters.
The price is about 80 encoder passes per photo. Negative results are reported as they are.
